\documentclass[border=10pt]{standalone}
\usepackage{fontspec}
\AtBeginDocument{\hyphenation{zenodo}}% never hyphenate "zenodo" inside a DOI
\setmainfont{Noto Serif}[Ligatures=TeXOff]
\XeTeXgenerateactualtext=1
\newfontfamily\bengalifont{Noto Serif Bengali}[Script=Bengali, NFSSFamily=bnfam]
\usepackage[Latin,Bengali]{ucharclasses}
% group-free font switching (as in header-common.tex): robust inside braces, cells and quotes
\makeatletter
\def\BN@saved{\rmdefault}
\def\BN@bn{bnfam}
\newcommand\BNin{\edef\BN@tmp{\f@family}\ifx\BN@tmp\BN@bn\else\let\BN@saved\BN@tmp\fi\fontfamily{bnfam}\selectfont}
\newcommand\BNout{\fontfamily{\BN@saved}\selectfont}
\makeatother
\setTransitionsFor{Bengali}{\BNin}{\BNout}
% Typographic quotes and dashes are not in a ucharclasses class (class 0), so a Bangla word right
% after “ or ‘ would not trigger the Bangla font. Treat them as Latin so that every switch fires.
\XeTeXcharclass"2018=\XeTeXcharclass`A
\XeTeXcharclass"2019=\XeTeXcharclass`A
\XeTeXcharclass"201C=\XeTeXcharclass`A
\XeTeXcharclass"201D=\XeTeXcharclass`A
\XeTeXcharclass"2013=\XeTeXcharclass`A
\XeTeXcharclass"2014=\XeTeXcharclass`A
\XeTeXcharclass"2026=\XeTeXcharclass`A
\usepackage{newunicodechar}
\newunicodechar{।}{{\bengalifont\char"0964}}
\newunicodechar{॥}{{\bengalifont\char"0965}}

\usepackage[table]{xcolor}
\usepackage{booktabs,tabularx,array}
\definecolor{bapilrule}{HTML}{1F2937}
\definecolor{bapilstripe}{HTML}{F3F4F6}
\renewcommand{\arraystretch}{1.28}
\setlength{\tabcolsep}{5pt}

\begin{document}
\begin{minipage}{17cm}
{\large\bfseries সারণি ৯।}\ {\large ব্যবহারকারী-চরিত্র (Persona)}\par\vspace{7pt}
\small
\rowcolors{2}{bapilstripe}{white}
\begin{tabularx}{\linewidth}{@{}>{\hsize=0.840\hsize\raggedright\arraybackslash}X>{\hsize=0.988\hsize\raggedright\arraybackslash}X>{\hsize=1.185\hsize\raggedright\arraybackslash}X>{\hsize=0.988\hsize\raggedright\arraybackslash}X@{}}
\arrayrulecolor{bapilrule}\toprule
\rowcolor{white}\textbf{চরিত্র} & \textbf{লক্ষ্য} & \textbf{বর্তমান সমস্যা} & \textbf{সবচেয়ে গুরুত্বপূর্ণ প্রয়োজনীয়তা} \\
\midrule
\textbf{ক. নীতি-বিশ্লেষক}, পরিকল্পনা অনুবিভাগ & নিয়ন্ত্রক পাঠ ও অতীত অঙ্গীকার দ্রুত খুঁজে পাওয়া; সারসংক্ষেপে উদ্ধৃত করা & বহু ওয়েবসাইটে খোঁজা; স্ক্যানকৃত PDF; কোন সংস্করণ চলমান তা অনিশ্চিত & FR-01, FR-02, FR-03 \\
\textbf{খ. বাজেট বিশ্লেষক বা সাংবাদিক} & বছর ও মন্ত্রণালয়ভেদে বরাদ্দ তুলনা & PDF-এ আটকে থাকা সারণি; বছরভেদে অসামঞ্জস্যপূর্ণ শিরোনাম & FR-04, NFR-07 \\
\textbf{গ. জন-পরামর্শ পরিচালনাকারী}, মন্ত্রণালয় & হাজারো মতামত ন্যায্যভাবে বোঝা & হাতে পড়া; সংখ্যালঘু মত হারিয়ে যাওয়ার ঝুঁকি & FR-05, GR-03 \\
\textbf{ঘ. গবেষক}, বিশ্ববিদ্যালয় & বাস্তব নীতি-টেক্সটে বাংলা পদ্ধতির বেঞ্চমার্ক & উন্মুক্ত ক্ষেত্রভিত্তিক বেঞ্চমার্ক বা নিয়ন্ত্রিত কর্পাস নেই & FR-08, NFR-06 \\
\textbf{ঙ. জ্যেষ্ঠ কর্মকর্তা} & সুপারিশের পেছনের প্রমাণে আস্থা & এআই-ফলাফলের উৎস বা সীমাবদ্ধতা দেখা যায় না & NFR-08, GR-01 \\
\textbf{চ. নাগরিক বা নাগরিক সমাজের পর্যবেক্ষক} & জনস্বার্থমূলক গবেষণায় এআই-এর ব্যবহার জানা & এআই ব্যবস্থার কোনো জনসমক্ষে নিবন্ধন নেই & FR-09, GR-06 \\
\bottomrule
\end{tabularx}
\par\vspace{3pt}{\scriptsize\color{gray} উৎস: BAPIL টেকনিক্যাল ব্লুপ্রিন্ট, সংস্করণ ১.০ (২০২৬)। CC BY 4.0।}
\end{minipage}
\end{document}
